\section{引言}

2-bit大语言模型量化把连续权重压入稀疏离散空间。向量量化以一个索引表示多个权重，能够利用码字shaping提高表示效率，也使优化变量成为沿注意力、MLP和残差流相互作用的整数assignment集合。AQLM、QuIP\#与QTIP说明表示几何、二阶敏感度和可实现编码需共同设计\cite{aqlm,quipsharp,qtip}；但部署对象仍是硬整数索引，soft loss或独立局部代理不会自动授权一组切换。

本文从block-scaled shared-codebook VQ出发。Scale先把不同动态范围的block映射到共同归一化领域，使整层共享4096个六维码字；Vector-GPTQ在top-128几何候选中加入激活Hessian代理与误差反馈，形成硬锚点。锚点之后，当前码字8-way邻域中的assignment是不增加码本或逻辑码率的剩余自由度，可由binary Concrete训练；真正困难不是开关能否获得梯度，而是最终硬动作是否同时保持任务与语言行为。

此前正式实验逐步排除了简单解释。聚合一阶在98.20\%向量上找到下降邻居，四视图严格共识把资格率降到48.71\%，局部曲率又优于共识排序，但硬集合仍失败。因果集合条件搜索随后在已接受状态下真实测量bundle边际收益：最后四层的四个坐标使训练目标改善2.73\%、held-out任务loss改善1.09\%，却让文本validation CE恶化0.919\%。这说明集合计分可被真实forward缓解，剩余瓶颈是task-only目标允许消费语言保持裕量。

V14提出唯一问题：若文本保持不是加权项，而是候选接受时不可交易的约束，固定局部动作空间中是否仍存在正任务收益方向？我们固定模型、source、28个curvature top-128 bundle、层序和数据；先测真实task gain，仅给改善候选计算固定top-32 teacher分布上的exact full-vocabulary text-train CE。候选必须满足CE不增加，才按任务loss排序。Validation和audit在搜索中不可见。零退化是检验“无需交易方向是否存在”的机制探针，不是现实部署的唯一容忍度；任何非零预算都属于新的预注册实验。

Figure~\ref{fig:v14}给出核心观察。28个bundle中27个降低任务目标，但所有27个都位于文本零退化边界之外；最小代价为$+0.00785\%$，四层均保持source。贡献为：（1）提出text-train CE先定义可行域、task gain再排序的词典序检验；（2）建立128-GiB prefix cache、候选batch与分块精确CE的单卡审计合同；（3）给出受限机制边界：当前固定动作集合中任务收益丰富而零文本代价方向缺失，后续应改变动作空间，而不是只对同一集合重排序。
